\chapter{从模形式到自守表示 III}

% ===== PDF p.033 =====
\textbf{上次}：
\begin{itemize}
\item 设 $\phi\in C^\infty(\GL_2(\Q)\backslash\GL_2(\A))$。
\item 则 $\phi=\phi_f$（某个 $f\in M_k$）当且仅当
\begin{itemize}
\item $\phi(g,g_\infty z)=z^{-k}\phi(g,g_\infty)$，对 $(g,g_\infty)\in\GL_2(\A)$、$z\in\C^\times$；
\item $\phi$ 适度增长；
\item 对每个 $g\in\GL_2(\Af)$ 有 $Y*\phi(g,-)=0$，其中 $Y\in\g_\C=\mathfrak{gl}_2(\R)_\C$ 是适当构造的自然元素。
\end{itemize}
\item $f\in M_k$ 是尖点形式，当且仅当对每个真抛物子群 $B\subseteq\GL_{2,\Q}$（即 Borel 子群），以 $N\subseteq\GL_{2,\Q}$ 记其幺幂根基，有
\[
\int_{N(\Q)\backslash N(\A)}\phi_f(ng)\,\mathrm{d}n=0
\]
对一切 $g\in\GL_2(\A)$ 成立。
\end{itemize}

\textbf{今天}：
\begin{itemize}
\item 把 $S_k$ 描述为 $\GL_2(\Af)$-表示。
\item 由此得到自守表示的主要例子来源。\znoteplain{原文脚注：更精确地说，$\widetilde\Phi(S_k)\subseteq L^2([\GL_2])$ 的闭包的不可约子表示，正是与尖点形式相伴的（$L^2$ 意义下的）自守表示。}
\end{itemize}

考察整个空间 $M_k$、$S_k$（而不只是某个紧开子群 $K\subseteq\GL_2(\Af)$ 对应的 $M_k(K)$、$S_k(K)$，也不只是同余子群 $\Gamma\subseteq\SL_2(\Z)$ 对应的 $M_k(\Gamma)$、$S_k(\Gamma)$）的好处是：现在可以使用局部 profinite 群 $\GL_2(\Af)$ 的表示论。就此而言，下面的观察很重要。

\textbf{观察}：
\begin{itemize}
\item $M_k$ 与 $S_k$ 是光滑且容许的 $\GL_2(\Af)$-表示（在下述意义下）。
\end{itemize}

\noindent\textbf{定义 5.1}\phantomsection\label{def:5.1}\quad 设 $G$ 为局部 profinite 群（例如 $\GL_2(\Af)$）。$G$ 在 $\C$-向量空间 $V$ 上的表示称为\textbf{光滑的}（smooth），若
\[
V=\bigcup_{K\subseteq G}V^K,
\]
其中 $K$ 遍历 $G$ 的紧开子群，且
\[
V^K:=\{v\in V\mid kv=v\ \text{对一切}\ k\in K\}.
\]
等价地，$V$ 光滑当且仅当作用态射 $G\times V\to V$ 连续，其中 $V$ 取离散拓扑。记 $\mathrm{Rep}^\infty_\C G$ 为 $G$ 在 $\C$-向量空间上的光滑表示构成的（abelian）范畴。

% ===== PDF p.034 =====
\noindent\textbf{定义 5.2}\phantomsection\label{def:5.2}\quad 设 $G$ 为局部 profinite 群。$G$ 的光滑表示 $V$ 称为\textbf{容许的}（admissible），若
\[
\dim_\C V^K<\infty
\]
对一切紧开子群 $K\subseteq G$ 成立。\znoteplain{光滑性与容许性是两件不同的事：光滑只要求每个向量都被某个紧开子群固定，容许性进一步要求每个 $V^K$ 有限维。对 $\GL_2(\Af)$ 而言 $M_k(K)$ 有限维，所以 $M_k$、$S_k$ 两者兼备；$C^\infty(G(\Q)\backslash G(\A))$ 光滑但不容许；$L^2([G])$ 连光滑都不是。}

看几个例子。
\begin{itemize}
\item $M_k=\bigcup_{K\subseteq\GL_2(\Af)}M_k^K$，其中 $M_k^K=M_k(K)$，它是光滑且容许的，因为各 $M_k(K)$ 都有限维。
\item 同理，$S_k$ 是 $\GL_2(\Af)$ 的容许表示。
\item 若 $G$ 在 $\Q$ 上约化，则 $C^\infty(G(\Q)\backslash G(\A))$ 是光滑的 $G(\Af)$-表示，但不容许。
\item $L^2([G])$ 不是光滑的 $G(\Af)$-表示。
\end{itemize}

在第 11 讲中，我们将对 $\Q$ 上的约化群 $G$ 引入自守形式空间。它可作为 $M_k$ 的替代（实际上是极大的推广），并且仍是光滑的 $G(\Af)$-表示。我们将会看到，光滑表示比 $L^2([G])$ 这样的酉表示在性质上更代数，这使得局部 profinite 群的光滑表示更便于操作。

在模形式的经典理论中，Hecke 算子（其定义看起来可能有些 ad hoc）很重要。从表示论的观点看，它们来自一个 Hecke 代数\idx{Hecke 代数}的作用。

\textbf{Hecke 代数}（局部 profinite 群的）：
\begin{itemize}
\item $G$ 为局部 profinite 群。
\begin{itemize}
\item 例如 $\GL_2(\Af)$、$\GL_n(\Q_p)$、$\Gal(\overline\Q/\Q)$、$\dots$
\end{itemize}
\item 对紧开子群 $K\subseteq G$，存在唯一的左不变（相应地，唯一的右不变）Haar 测度 $\mu$，使 $\mu(K)=1$。
\item 具体地：
\begin{itemize}
\item 由平移不变性，只需对紧开子群 $K'\subseteq G$ 定义 $\mu$。
\item 令 $K'':=K'\cap K$，它紧开。则 $K''$ 在 $K$ 中、相应地也在 $K'$ 中指数有限。
\item 令
\[
\mu(K'):=\frac{[K':K'']}{[K:K'']}.
\]
\item 这定义了所要的测度，且这是唯一可能的选择。
\item 注意 $\mu$ 的取值实际上在 $\Q$ 中。这是「局部 profinite 群的光滑表示论更代数」的第一个迹象。
\end{itemize}
\item 今后固定 $G$ 上一个左不变 Haar 测度 $\mu$。
\item 定义相应的「Hecke 代数」
\[
H(G):=C_c^\infty(G),
\]
即局部常值、紧支集函数 $G\to\C$ 之集。
\item $H(G)$ 通过卷积成为代数：
\[
f*g(x):=\int_G f(xy)g(y^{-1})\,\mathrm{d}y
\]
% ===== PDF p.035 =====
对 $f,g\in H(G)$。注意 $H(G)$ 隐含依赖于所选测度。不过可以把 $H(G)$ 内蕴地定义为紧支分布之空间，乘法由分布的卷积给出，参见 \texttt{[BRa, Definition 7]}。
\item $H(G)$ 中的乘法是纯代数运算：
\begin{itemize}
\item 每个 $f\in H(G)$ 都是有限和
\[
f=\sum_{i=1}^m a_i\chi_{g_iK_i},
\]
其中 $g_i\in G$，$K_i\subseteq G$ 为紧开子群。
\item 不妨设 $K_1=\dots=K_m$（通过缩小）。
\item 设 $K\subseteq G$ 紧开，$g_1,g_2\in G$。则存在 $h_j\in G$ 使
\[
g_1Kg_2K=\bigsqcup_{j=1}^m h_jK,\qquad
\chi_{g_1K}*\chi_{g_2K}=\sum_{j=1}^m\mu(K)\chi_{h_jK}.
\]
\end{itemize}
\item 代数 $H(G)$ 有很多幂等元，即
\[
e_K:=\frac{1}{\mu(K)}\chi_K,
\]
其中 $K\subseteq G$ 遍历紧开子群，$\chi_K$ 为 $K$ 的特征函数。但一般 $H(G)$ 没有单位元（除非 $G$ 紧）。
\end{itemize}

\textbf{Hecke 代数的模}（参见 \texttt{[BRa]}）：
\begin{itemize}
\item $H(G)$ 是群代数 $\C[G]$ 的替代品，它更适合用来研究 $G$ 的光滑表示。
\item 若 $V\in\mathrm{Rep}^\infty_\C G$，则
\[
f*v:=\int_G f(h)hv\,\mathrm{d}h
\]
（$f\in H(G)$，$v\in V$）定义了 $H(G)$ 在 $V$ 上的作用。
\item 具体地：若 $K\subseteq G$ 是紧开子群且充分小，使 $v\in V$ 被 $K$ 固定，而 $f=\sum_{i=1}^m a_i\chi_{g_iK}$（$a_i\in\C$，$g_i\in G$），则
\[
f*v=\sum_{i=1}^m a_ig_iv.
\]
\item 设 $K\subseteq G$ 为紧开子群。定义 $G$ 相对于 $K$ 的 Hecke 代数
\[
H(G,K):=e_KH(G)e_K.
\]
则 $H(G,K)$ 是以 $e_K$ 为单位的代数。
\item 设 $V$ 为 $G$ 的任意光滑表示。则不难看出
\[
V^K=e_K*V
\]
对任意紧开子群 $K\subseteq G$ 成立。
% ===== PDF p.036 =====
\item 把它应用于光滑 $G$-表示 $H(G)$（经左、右乘法），得到
\[
H(G,K)\cong C_c^\infty(K\backslash G/K),
\]
即 $H(G,K)$ 等同于 $G$ 上 $K$-双不变、紧支集函数之集。
\end{itemize}

下面这个小命题在随后讨论 $\GL_2(\Af)$-表示 $S_k$ 时很重要。
\begin{itemize}
\item 若 $V\in\mathrm{Rep}^\infty_\C G$ 不可约且 $V^K\neq0$，则 $V^K$ 作为 $H(G,K)$-模不可约。
\begin{itemize}
\item 事实上：若 $M\subseteq V^K$ 是非零 $H(G,K)$-子模，则由 $V$ 的不可约性有 $V=H(G)*M$，从而
\[
V^K=e_K(H(G)*M)=e_KH(G)e_KM=M,
\]
因为 $e_K*M=M$。
\end{itemize}
\item 反之可以证明：若对一切紧开子群 $K\subseteq G$，$V^K$ 不可约或为零，则 $V$ 不可约（当 $V\neq0$ 时）。
\end{itemize}

也就是说，这个性质使我们能把 Hecke 特征系统\idx{Hecke 特征系统}联系到不可约容许的 $\GL_2(\Af)$-表示上。

\textbf{从不可约容许 $\GL_2(\Af)$-表示到 Hecke 特征系统}：
\begin{itemize}
\item 设 $V$ 为不可约、容许的 $\GL_2(\Af)$-表示。
\item 回忆紧开子群
\[
K(m)=\ker(\GL_2(\widehat\Z)\to\GL_2(\Z/m)),\qquad m\ge0,
\]
构成 $\GL_2(\Af)$ 中单位元的紧开邻域基。
\item 由上述重要性质与容许性（后者容易推出 Schur 引理的一个版本）知
\begin{itemize}
\item 当 $m\gg0$ 时，$V^{K(m)}$ 是不可约 $H(\GL_2(\Af),K(m))$-模；
\item $\End_{H(\GL_2(\Af),K(m))}(V^{K(m)})\cong\C$。
\end{itemize}
\item 下面我们用到：对每个素数 $p$，光滑 $\GL_2(\Af)$-表示沿包含 $\GL_2(\Q_p)\subseteq\GL_2(\Af)$ 限制，得到光滑的 $\GL_2(\Q_p)$-表示。
\item 令
\[
S:=\bigcup_{m\ge0,\,V^{K(m)}\neq0}\{\,p\ \text{素数}:p\mid m\,\}.
\]
\item 注意 $p\in S$ 当且仅当 $V^{\GL_2(\Z_p)}=0$。
\item 若 $p\notin S$，即 $V^{\GL_2(\Z_p)}\neq0$，则称素数 $p$ 对 $V$ \textbf{非分歧\idx{非分歧}}。
\item 下面的观察非常重要：对每个素数 $p$，代数 $H(\GL_2(\Q_p),\GL_2(\Z_p))$ 交换！
\item 更强的是：取 $m\ge1$ 与 $p\nmid m$。则 $H(\GL_2(\Q_p),\GL_2(\Z_p))$ 在 $H(\GL_2(\Af),K(m))$ 中中心。更精确地说，对某个 $A\in\GL_2(\Q_p)$ 的双陪集，可以联系到元素
\[
(1,\dots,1,A,1,\dots,1)\in\GL_2(\Af)
\]
（第 $p$ 个分量为 $A$，其余分量为单位元）的双陪集，而所得的嵌入
\[
H(\GL_2(\Q_p),\GL_2(\Z_p))\longrightarrow H(\GL_2(\Af),K(m))
\]
% ===== PDF p.037 =====
其像落在 $H(\GL_2(\Af),K(m))$ 的中心。因为容易看出：由 $(1,\dots,1,A,1,\dots,1)$（$A\in\GL_2(\Q_p)$）给出的双陪集与 $B\in\GL_2(\Af)$ 的双陪集交换，只要 $B$ 在 $p$ 处的分量是单位矩阵。于是 $H(\GL_2(\Q_p),\GL_2(\Z_p))$ 的中心性实际上由它的交换性推出。
\item 为证交换性，我们用 Gelfand 技巧。即考虑反对合
\[
\sigma:\GL_2(\Q_p)\to\GL_2(\Q_p),\qquad g\mapsto g^{tr},
\]
它具有以下性质：
\begin{itemize}
\item $\sigma$ 保持 $\GL_2(\widehat\Z_p)$；
\item $\sigma(f_1*f_2)=\sigma(f_2)*\sigma(f_1)$，对 $f_1,f_2\in H(\GL_2(\Q_p),\GL_2(\Z_p))$；
\item 双陪集 $\GL_2(\Z_p)g\GL_2(\Z_p)$（$g\in\GL_2(\Q_p)$ 为对角矩阵）张成 $H(\GL_2(\Q_p),\GL_2(\Z_p))$（由初等因子定理）；
\item 现在可以证明交换性了。因为 $\sigma(f)=f$ 对一切 $f\in H(\GL_2(\Q_p),\GL_2(\Z_p))$ 成立，故
\[
f_1*f_2=\sigma(f_1*f_2)=\sigma(f_2)*\sigma(f_1)=f_2*f_1,
\]
如所欲。
\end{itemize}
\item $H(\GL_2(\Q_p),\GL_2(\Z_p))$ 的中心性很重要，因为它允许定义 $V$ 的 Hecke 特征系统（回忆 $V$ 是不可约容许的 $\GL_2(\Af)$-表示）。
\item 回忆 $S$ 记 $V$ 的分歧素数集，即那些使 $V^{\GL_2(\Z_p)}=0$ 的素数 $p$。
\item 对 $p\notin S$，我们下面构造一个 $\C$-代数态射
\[
\sigma_{V,p}:H(\GL_2(\Q_p),\GL_2(\Z_p))\to\C,
\]
它自然联系于 $V$。
\item 为此取 $p\notin S$，并选 $v\in V$ 被 $\GL_2(\Z_p)$ 固定。则由 $\GL_2(\Af)$-作用的光滑性，$v$ 被某个 $K(m)$ 固定，其中 $m\ge0$ 且 $p\nmid m$。
\item 如前所述，$V$ 的不可约性蕴涵 $V^{K(m)}$ 不可约，从而由 Schur 引理
\[
\End_{H(\GL_2(\Af),K(m))}(V^{K(m)})\cong\C.
\]
\item 由于 $H(\GL_2(\Q_p),\GL_2(\Z_p))$ 在 $H(\GL_2(\Af),K(m))$ 中中心，它通过 $H(\GL_2(\Af),K(m))$-模同态作用在 $V^{K(m)}$ 上。
\item 于是得到态射
\[
\sigma_{V,p}:H(\GL_2(\Q_p),\GL_2(\Z_p))\to\End_{H(\GL_2(\Af),K(m))}(V^{K(m)})\cong\C,
\]
它先验依赖于 $v$ 与 $m$。但不难看出 $\sigma_{V,p}$ 在放大 $m$ 时不变。这意味着 $\sigma_{V,p}$ 实际上与 $v$、$m$ 无关，从而典范地联系于 $V$ 与 $p$。
\end{itemize}

\noindent\textbf{定义 5.3}\phantomsection\label{def:5.3}\quad 系统
\[
\{\sigma_{V,p}:H(\GL_2(\Q_p),\GL_2(\Z_p))\to\C\}_{p\notin S}
\]
称为 $V$ 的\textbf{Hecke 特征系统}（system of Hecke eigenvalues）。

% ===== PDF p.038 =====
我们给出 Hecke 特征系统的另一种视角。设 $S$ 为某个有限素数集，定义「远离 $S$ 的 Hecke 代数」为
\[
T_S:=H(\GL_2(\Af^S),\GL_2(\widehat\Z^S)),
\]
其中
\[
\widehat\Z^S:=\prod_{p\notin S}\Z_p
\qquad\text{且}\qquad
\Af^S:=\prod_{p\notin S}\Q_p.
\]
\znote{原文把 $\Af^S$ 写成 $\Q\otimes_\Z\widehat\Z^S$；受限张量积的正确刻画是 $\Af^S=\prod_{p\notin S}\Q_p$（即几乎处处落在 $\Z_p$ 中的族），张量积一般不保持无穷乘积，故此处按正解译。[来源：原文 p.038 紧随其后的 $\widehat\Z^S$ 定义与 $T_S$ 的受限张量积式，参见 \texttt{[Fla79]}。]}

利用 Gelfand 技巧可以看出 $T_S$ 交换。此外，若展开适当的定义（参见 \texttt{[Fla79]}），则「受限张量积」给出
\[
T_S\cong\sideset{}{'}\bigotimes_{p\notin S}H(\GL_2(\Q_p),\GL_2(\Z_p)).
\]

现在还可以作如下定义。

\noindent\textbf{定义 5.4}\phantomsection\label{def:5.4}\quad \textbf{Hecke 特征系统}（Hecke eigensystem，或 system of Hecke eigenvalues）是 $T_S$ 中的一个极大理想，其中 $S$ 为有限素数集，$T_S:=H(\GL_2(\Af^S),\GL_2(\widehat\Z^S))$。

我们可以收集以下观察。
\begin{itemize}
\item $T_S$ 的维数可数，因此每个 Hecke 特征系统的剩余域为 $\C$。
\item $S$ 的一个 Hecke 特征系统等价于一族 $\C$-代数同态
\[
\sigma_p:H(\GL_2(\Q_p),\GL_2(\Z_p))\to\C\qquad(p\notin S),
\]
如前。
\item 若 $S\subseteq S'$ 为两个有限素数集，则每个 $S$ 的 Hecke 特征系统诱导 $S'$ 的一个。两个 Hecke 特征系统 $\{\sigma_p\}_{p\notin S}$ 与 $\{\sigma'_p\}_{p\notin S'}$ 称为\textbf{等价的}，若它们在某个包含 $S$ 与 $S'$ 的有限素数集之外一致。
\item 如前所见，每个不可约容许 $\GL_2(\Af)$-表示都通过上述构造给出一个 Hecke 特征系统。
\end{itemize}

我们将会看到，在 Langlands 纲领中，与模形式相伴的 Hecke 特征系统被期望具有算术意义。

为了对 Hecke 特征系统获得更有用的视角，我们需要 $H(\GL_2(\Q_p),\GL_2(\Z_p))$ 的完整描述。

\textbf{$H(\GL_2(\Q_p),\GL_2(\Z_p))$ 的完整描述}：
\begin{itemize}
\item 暂令 $G:=\GL_2(\Q_p)$，$K=\GL_2(\Z_p)$。
\item 归一化 $G$ 上的 Haar 测度，使 $\mu(K)=1$。
\item 由初等因子定理，$\C$-向量空间 $H(G,K)$ 以
\[
K\begin{pmatrix}p^i&0\\0&p^j\end{pmatrix}K,\qquad
(i,j)\in\Z^2_+:=\{(i,j)\in\Z^2\mid i\ge j\}
\]
为基。
% ===== PDF p.039 =====
\item 我们断言
\[
f:\C[X_1,X_2^{\pm1}]\xrightarrow{\ \sim\ }H(G,K)
\]
由
\[
X_1\mapsto K\begin{pmatrix}p&0\\0&1\end{pmatrix}K,\qquad
X_2\mapsto K\begin{pmatrix}p&0\\0&p\end{pmatrix}K
\]
给出。
\item 这里（以及今后）我们把双陪集与其特征函数等同。
\item 我们来验证这个断言（不借助 Gelfand 技巧）：
\begin{itemize}
\item 元素 $K\begin{pmatrix}p&0\\0&p\end{pmatrix}K$ 在 $H(G,K)$ 中中心，因为对一切 $g\in G$ 与群论中心中的 $z\in G$ 有
\[
KzK\cdot KgK=KgK\cdot KzK.
\]
\item 这说明 $f$ 良定义（注意 $\C[X_1,X_2^{\pm1}]$ 交换，而 $H(\GL_2(\Q_p),\GL_2(\Z_p))$ 先验不交换）。
\item 再验证 $f$ 满射。乘以 $f(X_2)$ 之后，只需说明每个
\[
K\begin{pmatrix}p^i&0\\0&1\end{pmatrix}K\qquad(i\ge1)
\]
都在 $f$ 的像中。
\item 用归纳法可见：对 $n\ge1$，$n$ 重乘积
\[
K\begin{pmatrix}p&0\\0&1\end{pmatrix}K\cdot\ldots\cdot
K\begin{pmatrix}p&0\\0&1\end{pmatrix}K
\]
与 $\mathrm{Mat}_{2,2}(\Z_p)$ 中行列式赋值等于 $n$ 的矩阵之集 $\mathrm{Mat}_{2,2}(\Z_p)^{\mathrm{val}=n}$ 一致。
\item 由 $H(G,K)$ 中乘法的定义可以推出，上述乘积就是 $\mathrm{Mat}_{2,2}(\Z_p)^{\mathrm{val}=n}$ 的特征函数 $\chi$。
\item 由于这个特征函数 $\chi$ 以 $K\begin{pmatrix}p^n&0\\0&1\end{pmatrix}K$ 的特征函数为系数 $1$ 的求和项，我们可以（用一个小归纳）断定它落在 $f$ 的像中，因为 $\mathrm{Mat}_{2,2}(\Z_p)^{\mathrm{val}=n}$ 中出现的其它双陪集都有形如
\[
\begin{pmatrix}p^j&0\\0&p^i\end{pmatrix},\qquad i,j\ge1,\ i+j=n
\]
的矩阵作为代表元。
\item 特别地，我们看到 $H(G,K)$ 必定交换，而无需借助 Gelfand 技巧。
% ===== PDF p.040 =====
\item 现在不难断定 $f$ 必为同构。事实上，
\[
\Spec(H(G,K))\subseteq\Spec(\C[X_1,X_2^{\pm1}])
\]
定义了一个子概形，它到 $\Spec(\C[X_2^{\pm1}])$ 的投影是无限秩自由的；而这对 $\Spec(\C[X_1,X_2^{\pm1}])$ 自身才可能（利用后者的整性）。
\end{itemize}
\end{itemize}

我们总结关于 Hecke 特征系统的讨论。

\textbf{要点}：
\begin{itemize}
\item 每个不可约容许 $\GL_2(\Af)$-表示 $V$ 都典范地给出一个 Hecke 特征系统
\[
\tilde a_p\in\C,\qquad \tilde b_p\in\C^\times
\]
对 $V$ 的分歧素数有限集 $S$ 之外的一切素数 $p$。
\item 具体地：若 $v\in V$ 被 $K(m)$ 固定且 $p\nmid m$，则
\[
\tilde T_p(v):=\GL_2(\Z_p)\begin{pmatrix}p&0\\0&1\end{pmatrix}\GL_2(\Z_p)*v=\tilde a_pv,
\]
其中 $\tilde T_p$ 是「$p$ 处的 Hecke 算子」，且
\[
\tilde S_p(v):=\GL_2(\Z_p)\begin{pmatrix}p&0\\0&p\end{pmatrix}\GL_2(\Z_p)*v=\tilde b_pv.
\]
\item 注意 Hecke 特征系统完全由复数族 $\{\tilde a_p,\tilde b_p\}_{p\notin S}$ 决定。
\item 令 $Z\subseteq\GL_2$ 为中心，即标量矩阵构成的子群。
\item 令 $\psi:\Af^\times\cong Z(\Af)\to\C^\times$ 为 $V$ 的中心特征（由 Schur 引理存在）。则
\[
\tilde b_p=\psi((1,\dots,1,p,1,\dots,1))
\]
对满足 $Z_p^\times\subseteq\Af^\times$ 在 $V$ 上平凡作用的素数 $p$ 成立。
\item 中心特征 $\psi$ 给出不可约容许 $\GL_2(\Af)$-表示的另一个不变量。
\item 注意 $\psi$ 一般不由 Hecke 特征系统决定，因为缺少有限多个 $p$，且没有要求 $V$ 在 $\Q^\times$ 上平凡，即没有要求「自守性」。
\end{itemize}

我们现在用上述讨论得到光滑容许 $\GL_2(\Af)$-表示 $S_k$ 的分解。先收集一些一般性观察，再得到精确的分解。

\textbf{关于 $\GL_2(\Af)$-表示 $S_k$ 的一般结果}：
\begin{itemize}
\item $\GL_2(\Af)$-表示 $S_k$ 分解为不可约容许表示 $V_i$ 的直和
\[
S_k\cong\bigoplus_{i\in I}V_i.
\]
\item 事实上：
% ===== PDF p.041 =====
\begin{itemize}
\item 在乘以特征 $|\det|_{\mathrm{ad\acute elique}}^{k/2}|_{\GL_2(\Af)}$ 的意义下，有一个 $\GL_2(\Af)$-等变嵌入 $S_k\subseteq L^2([\GL_2])$，从而 $\C$-向量空间 $S_k\otimes|\det|_{\mathrm{ad\acute elique}}^{k/2}$ 带有非退化的 $\GL_2(\Af)$-等变 Hermite 配对
\[
(\cdot,\cdot):S_k\times S_k\to\C,\qquad
(f,g)\mapsto\int_{[\GL_2]}\tilde\phi_f\overline{\tilde\phi_g}.
\]
\item 由此，通过取正交补，每个不可约 $\GL_2(\Af)$-子商其实已经是直和项。
\item 由 Zorn 引理，每个非零光滑 $\GL_2(\Af)$-表示 $V$ 都有不可约子商。事实上，不妨设 $V$ 由某个 $v\in V$ 生成；再对不含 $v$ 的子模之集合应用 Zorn 引理。
\end{itemize}
\end{itemize}

我们需要 Shalika 与 Piatetski-Shapiro 的下述结果。

\noindent\textbf{定理 5.5}（Shalika/Piatetski-Shapiro）\phantomsection\label{thm:5.5}\quad 对每个有限素数集 $S$ 与每个 Hecke 特征系统 $\{\tilde a_p,\tilde b_p\}_{p\notin S}$，至多存在一个具有该（等价类）Hecke 特征系统的不可约子表示 $V\subseteq S_k$。

\begin{itemize}
\item 这是 $\GL_n$ 尖点自守表示的「强重数一定理」的特例，参见 \texttt{[GH19, Theorem 11.7.2]}。\znoteplain{原文脚注：更精确地说，我们通过尖点自守形式的稠密性 \texttt{[GH19, Theorem 6.5.1.]} 与嵌入 $S_k\otimes|\det|_{\mathrm{ad\acute elique}}^{k/2}$，从 $L^2_{\mathrm{cusp}}([\GL_2])$ 的上述更强结果推出它。}
\item 我们可以记录下述推论：设 $0\neq f\in S_k$ 被 $K(m)$ 固定，并且是 $p\nmid m$ 时 $\tilde T_p$、$\tilde S_p$ 算子的本征向量。则由 $f$ 生成的 $\GL_2(\Af)$-表示 $V:=\langle f\rangle_{\GL_2(\Af)}$ 不可约。\znoteplain{「强重数一」这个名字容易让人误解：它并不是说每个不可约表示的重数为 $1$，而是说给定的 Hecke 特征系统（等价类）至多对应一个不可约子表示。正因如此，下面的定理 5.6 才能把 $S_k$ 的分解按新形式逐个列出，而不出现重数。}
\begin{itemize}
\item 因为每个非零 $v\in V$ 都给出一个 Hecke 特征系统，它与 $f$ 的等价；由强重数一\idx{强重数一}，它们必须落在同一子空间中。
\end{itemize}
\end{itemize}

现在我们把 Hecke 算子 $\tilde T_p$、$\tilde S_p$ 与它们在模形式上的经典对应物联系起来，并通过 Atkin/Lehner 的新形式\idx{新形式}描述 $S_k$ 的分解。

\textbf{以新形式为指标的 $S_k$ 分解}：
\begin{itemize}
\item 由 Atkin--Lehner/Casselman 的一个结果（参见 \texttt{[Cas73]}），对每个不可约 $\GL_2(\Af)$-子表示 $V\subseteq S_k$，存在某个 $m\in\N$ 与某个 $0\neq v\in V$，使 $v$ 被子群
\[
K_1(m):=\Bigl\{\begin{pmatrix}a&b\\c&d\end{pmatrix}\Bigm|c\equiv0,\ d\equiv1\ (\mathrm{mod}\ m)\Bigr\}\subseteq\GL_2(\widehat\Z)
\]
固定。
\item 令 $\Gamma_1(m):=\GL_2(\Z)\cap K_1(m)$。
\item 注记：
\begin{itemize}
% ===== PDF p.042 =====
\item 我们的记号与许多文献不同，那里 $\Gamma_1(m)\subseteq\SL_2(\Z)$。这一记号冲突是无害的，因为
\[
\Gamma_1(m)\backslash\HH^\pm\cong(\Gamma_1(m)\cap\SL_2(\Z))\backslash\HH.
\]
\item 使这样的非零 $v\in V$ 存在的最小 $m$ 称为 $V$ 的\textbf{导子}（conductor）。此时 $v$ 在相差一个标量的意义下唯一（称为该表示的新形式）。\znoteplain{导子是「级」的表示论版本：它记录表示在哪些素数处真正分歧，而不是记录一个具体的同余子群。Atkin--Lehner 理论的内容正是：每个不可约子表示都恰有一个极小级，且该级上的新形式（相差标量）唯一；其余同级的本征形式都由更低级的形式「旧形式」诱导而来。}
\item 更多细节见 \texttt{[Del73]}。
\end{itemize}
\item 我们已经注意到，态射
\[
\Gamma_1(m)\backslash\HH^\pm\xrightarrow{\ \sim\ }\GL_2(\Q)\backslash(\GL_2(\Af)/K_1(m)\times\HH^\pm),\qquad
[z]\mapsto[(1,z)]
\]
在第 3 讲中是同构。当然，该同构与 $\omega^{\otimes k}$ 相容。
\item 特别地，
\[
M_k(\Gamma_1(m))\cong M_k(K_1(m)),
\]
并且两者都与 $\Gamma_1(m)\cap\SL_2(\Z)$ 的经典模形式空间一致。
\item 让我们把这个同构写得更明确，取 $\tilde f\in M_k(K_1(m))$。
\item 设 $f\in M_k(\Gamma_1(m))$ 是对应于 $\tilde f$ 的截面。
\item 则
\[
\tilde f(g,z)=\tilde f(\gamma,z)=j(\gamma^{-1},z)^{-k}f(\gamma^{-1}z),
\]
其中 $g=\gamma h$，$\gamma\in\GL_2(\Q)$，$h\in K_1(m)$。
\item 注意这里用到了 $\GL_2(\Q)K_1(m)=\GL_2(\Af)$。
\item 固定某个满足 $p\nmid m$ 的素数 $p$。
\item 两个算子 $\tilde T_p$、$\tilde S_p$ 作用在 $M_k(K_1(m))$ 上。于是由结构的搬移，它们作用在 $M_k(\Gamma_1(m))$ 上。
\item 我们把 $\tilde S_p$ 在 $M_k(\Gamma_1(m))$ 上的这个作用写明确。
\item 对矩阵 $\gamma\in\GL_2(\Q)$，记 $\gamma_p\in\GL_2(\Af)$ 为第 $p$ 个分量为 $\gamma\in\GL_2(\Q)\subseteq\GL_2(\Q_p)$、其余分量为单位矩阵的元素。不要把此元素与 $\gamma_{\mathrm{diag}}\in\GL_2(\Af)$ 混淆，后者指由对角嵌入 $\GL_2(\Q)\subseteq\GL_2(\Af)$ 得到的元素。
\item $\tilde S_p$ 在 $M_k(K_1(m))$ 上的作用容易写出，但要在 $M_k(\Gamma_1(m))$ 上描述它，需找表达式
\[
\begin{pmatrix}p&0\\0&p\end{pmatrix}_p=\gamma_{\mathrm{diag}}h,\qquad
\gamma\in\GL_2(\Q),\ h\in K_1(m).
\]
\item 因为 $p$、$m$ 互素，存在 $a,b\in\Z$ 使
\[
A:=\begin{pmatrix}a&b\\m&p\end{pmatrix}\in\GL_2(\Z).
\]
\item 则
\[
A_{\mathrm{diag}}\cdot\begin{pmatrix}p^{-1}&0\\0&p^{-1}\end{pmatrix}_{\mathrm{diag}}\cdot
\begin{pmatrix}p&0\\0&p\end{pmatrix}_p\in K_1(m),
\]
% ===== PDF p.043 =====
其中用了前面引入的记号 $(-)_{\mathrm{diag}}$、$(-)_p$。
\item 换言之，可以取
\[
\gamma:=\begin{pmatrix}p&0\\0&p\end{pmatrix}A^{-1}.
\]
\item 取此 $\gamma$，算子 $\tilde S_p$ 通过
\[
f\mapsto\bigl(z\mapsto j(\gamma^{-1},z)^{-k}f(\gamma^{-1}z)\bigr)
\]
作用在 $M_k(\Gamma_1(m))$ 上。
\item 群
\[
\Gamma_0(m):=\Bigl\{\begin{pmatrix}a&b\\c&d\end{pmatrix}\Bigm|c\equiv0\ (\mathrm{mod}\ m)\Bigr\}
\]
正规化 $\Gamma_1(m)$，于是有 $\Gamma_0(m)/\Gamma_1(m)\cong(\Z/m)^\times$ 的作用：
\[
(B,f)\mapsto j(B,z)^{-k}f(Bz),\qquad B\in\Gamma_0(m),\ f\in M_k(\Gamma_1(m)).
\]
\znote{原文把 $\Gamma_0(m)$ 写成 $c\equiv1\ (\mathrm{mod}\ m)$，应为 $c\equiv0$：只有 $c\equiv0$ 时 $\Gamma_0(m)$ 才包含 $\Gamma_1(m)$（后者要求 $c\equiv0$、$d\equiv1$），$\Gamma_0(m)/\Gamma_1(m)\cong(\Z/m)^\times$ 与旁特征 $\chi(d)$ 的作用才自洽。此处按 $c\equiv0$ 译。[来源：原文 p.042 的 $\Gamma_1(m)$ 定义与 p.043 旁特征作用的相容性；同一笔误亦见于 p.032。]}
\item 于是得到分解
\[
M_k(\Gamma_1(m))=\bigoplus_{\chi:(\Z/m)^\times\to\C^\times}M_k(\Gamma_0(m),\chi),
\]
其中 $M_k(\Gamma_0(m),\chi)$ 按定义是 $\Gamma_0(m)$ 上旁特征为 $\chi$ 的模形式空间，即满足
\[
f\Bigl(\frac{az+b}{cz+d}\Bigr)=\chi(d)(cz+d)^kf(z)
\]
对 $\begin{pmatrix}a&b\\c&d\end{pmatrix}\in\Gamma_0(m)$ 成立的 $\Gamma_1(m)$-模形式之空间。
\item 对每个 Dirichlet 特征 $\chi:(\Z/m)^\times\to\C^\times$，存在唯一的 adélic 特征
\[
\psi_\chi:\Q^\times\backslash\A^\times\bigl/((1+m\widehat\Z)\cap\widehat\Z^\times\bigr)\to\C^\times,
\]
使 $\psi_\chi(r)=r^{-k}$ 对 $r\in\R_{>0}$ 成立。
\item 我们得到
\[
M_k(\Gamma_0(m),\chi)\cong M_k(K_1(m),\psi_\chi),
\]
其中右端指 $\Af^\times$（$\GL_2(\Af)$ 的中心）在 $M_k(K_1(m))$ 上作用的 $\psi_\chi$-本征空间。
\item 用这套术语可以完成对算子 $\tilde S_p$ 在 $M_k(\Gamma_1(m))$ 上的描述。
\item 即：$\tilde S_p$ 在子空间 $M_k(\Gamma_0(m),\chi)\subseteq M_k(\Gamma_1(m))$ 上的作用，由
\[
\gamma^{-1}=\begin{pmatrix}p^{-1}&0\\0&p^{-1}\end{pmatrix}A
\]
的作用给出；而 $\begin{pmatrix}p^{-1}&0\\0&p^{-1}\end{pmatrix}$ 作用为 $p^k$，$A\in\Gamma_0(m)$ 作用为 $\chi(p)$。
\item 现在描述 $\tilde T_p$ 在 $M_k(\Gamma_0(m),\chi)$ 上的作用（由结构的搬移）。
\item 有
\[
\GL_2(\Z_p)\begin{pmatrix}p&0\\0&1\end{pmatrix}\GL_2(\Z_p)
=\bigsqcup_{j=0}^{p-1}\begin{pmatrix}p&j\\0&1\end{pmatrix}\GL_2(\Z_p)
\sqcup\begin{pmatrix}1&0\\0&p\end{pmatrix}\GL_2(\Z_p).
\]
% ===== PDF p.044 =====
\item 我们同样要把（上述代表的）双陪集分解为
\[
\gamma\cdot h,\qquad \gamma\in\GL_2(\Q),\ h\in K_1(m).
\]
\item 可以取分解
\[
\begin{pmatrix}p&j\\0&1\end{pmatrix}_p
=\begin{pmatrix}p&j\\0&1\end{pmatrix}_{\mathrm{diag}}\cdot h
\]
对 $j=0,\dots,p-1$，其中 $h\in K_1(m)$，即
\[
\gamma^{-1}=\begin{pmatrix}p^{-1}&-p^{-1}j\\0&1\end{pmatrix}.
\]
\item 对 $\begin{pmatrix}1&0\\0&p\end{pmatrix}$，与上面的论证类似，可得
\[
\gamma^{-1}=\begin{pmatrix}a&b\\m&p\end{pmatrix}\cdot\begin{pmatrix}1&0\\0&p^{-1}\end{pmatrix}
=\begin{pmatrix}1&0\\0&p^{-1}\end{pmatrix}\cdot\begin{pmatrix}a&p^{-1}b\\pm&p\end{pmatrix},
\]
如所欲。
\item 把这些放在一起，得到
\[
\tilde T_p(f)=p^k\chi(p)f(pz)+\sum_{j=0}^{p-1}f\Bigl(\frac{z-j}{p}\Bigr).
\]
\item 若把带旁特征的模形式 $f\in M_k(\Gamma_0(m),\chi)$ 写成 Fourier 展开
\[
f(z)=\sum_{n=0}^\infty a_nq^n,\qquad q=e^{2\pi iz},
\]
则更具体地
\[
\tilde T_p(f)(z)=\sum_{n=0}^\infty pa_{pn}q^n+\sum_{n=0}^\infty p^k\chi(p)a_nq^{pn}.
\]
\item 这由下述事实推出：
\[
\sum_{j=0}^{p-1}e^{2\pi i\frac{-jn}{p}}=
\begin{cases}
p,&p\mid n,\\
0,&\text{否则}.
\end{cases}
\]
\item 特别地，若 $f\in M_k(\Gamma_0(m),\chi)$ 是 $\tilde T_p$ 的本征向量且 $a_1=1$，则（考察 $q^1$ 前的系数）
\[
\tilde T_p(f)=pa_pf.
\]
这促使我们定义
\[
T_p=\frac1p\tilde T_p,
\]
它就是作用在模形式上的经典 Hecke 算子（如同 \texttt{[DS05, Section 5.2.]}）。换言之，$T_p$ 算子的 Hecke 特征值就是归一化本征形式的 Fourier 系数。\znoteplain{这里必须留意两套归一化：$\tilde T_p$ 是双陪集的特征函数（不除体积）给出的算子，而经典 $T_p$ 是它在 $\GL_2(\Af)$ 侧的「平均」版本，两者相差因子 $p$（正好是 $\tilde T_p$ 中 $p$ 个斜陪集的个数）。因此同一个本征形式满足 $\tilde T_p f=pa_pf$ 与 $T_pf=a_pf$，本征值 $\tilde a_p=pa_p$ 与 $a_p$ 只在差一个 $p$ 的意义下相同；比较两种文献时若忘了这个因子，就会得到相差 $p$ 的「矛盾」。}
\end{itemize}

% ===== PDF p.045 =====
我们现在可以证明关于 $\GL_2(\Af)$-表示 $S_k$ 分解的定理。

\noindent\textbf{定理 5.6}\phantomsection\label{thm:5.6}\quad $\GL_2(\Af)$-表示 $S_k$ 分解为不可约表示
\[
S_k\cong\bigoplus_f\langle f\rangle_{\GL_2(\Af)},
\]
其中 $f$ 遍历 $\Gamma_0(N)$ 上、旁特征为 $\chi$ 的（归一化）新形式，$N$ 与 Dirichlet 特征 $\chi:(\Z/N)^\times\to\C^\times$ 都在变。

我们不给出新形式的精确定义，参见 \texttt{[DS05, Section 5.8.]}。粗略地说，新形式是在所有具有等价 Hecke 特征系统的本征形式中实现极小 $N$ 的本征形式。

\textbf{证明.} 经过上述准备，我们现在可以引用 \texttt{[DS05, Proposition 5.8.4]}、\texttt{[DS05, Theorem 5.8.2]}，它们蕴涵：对 $S_k$ 中出现的每个 Hecke 特征系统，存在唯一的归一化（即 Fourier 展开中 $a_1=1$）新形式，其 Hecke 特征系统与之等价。\hfill$\square$

在其它教材中，例如 \texttt{[Gel75, Proposition 3.1]}，自守表示直接联系到某个 $f\in M_k(\Gamma_0(m),\chi)$，即不引入空间 $M_k$，而用 $\GL_2(\Af)$-等变嵌入
\[
M_k\to C^\infty(\GL_2(\Q)\backslash\GL_2(\A)),\qquad
f\mapsto\phi_f(g,g_\infty)=j(g_\infty,i)^{-k}f(g,g_\infty i),
\]
以及
\[
S_k\otimes|\det|_{\mathrm{ad\acute elique}}^{k/2}\to L^2([\GL_2]),\qquad
f\mapsto\tilde\phi_f(g,g_\infty)=|\det(g,g_\infty)|_{\mathrm{ad\acute elique}}^{k/2}\phi_f.
\]
我们指出，在上述同构 $M_k(\Gamma_1(m))\cong M_k(K_1(m))$ 下，我们的公式特殊化为 \texttt{[Gel75, Proposition 3.1.]} 中的公式。此外，\texttt{[Gel75, Lemma 3.7.]} 中出现的看似意外的因子 $p^{k/2-1}$ 有一个简单的解释：出现了 $|\det|_{\mathrm{ad\acute elique}}^{k/2}$ 的扭，以及从（我们的）$\tilde T_p$ 到经典 Hecke 算子 $T_p=\frac1p\tilde T_p$ 的过渡。
